%! Polynomial Functions and End Behavior — Unit 2, Lesson 2.5 — Homework (Answer Key)
\documentclass[10pt]{article}
\usepackage{precalculus-article}
\usepackage{precalculus-key}   % pulls in -boxes; do NOT also load -boxes
\newcommand{\TallMath}[1]{$\displaystyle #1\rule[-1.4em]{0pt}{3.2em}$}

\begin{document}

\pageheader{Unit 2, Lesson 2.5}{Homework --- Answer Key}

\vspace{0.12in}

\begin{enumerate}[itemsep=10pt]

\item From standard form. Name the leading term, then write both limit
      statements. Watch the exponents, not the order of the terms.

\smallskip
\renewcommand{\arraystretch}{1.6}
\begin{center}
\begin{tabular}{|l|c|c|c|}
\hline
\textbf{Polynomial} & \textbf{Leading term} & $\lim_{x \to \infty} p(x)$ & $\lim_{x \to -\infty} p(x)$ \\
\hline
$p(x) = 5x^4 - 3x + 8$      & \ans{$5x^4$} & \ans{$\infty$} & \ans{$\infty$} \\
\hline
$q(x) = -x^3 + 2x^2 - 9$    & \ans{$-x^3$} & \ans{$-\infty$} & \ans{$\infty$} \\
\hline
$r(x) = 6 - 2x^2 + x^5$     & \ans{$x^5$} & \ans{$\infty$} & \ans{$-\infty$} \\
\hline
$s(x) = -4x^{10} + x^{7}$   & \ans{$-4x^{10}$} & \ans{$-\infty$} & \ans{$-\infty$} \\
\hline
\end{tabular}
\end{center}

\item From factored form. Add the exponents for the degree, multiply the
      leading pieces for the leading term --- do not expand.

\begin{enumerate}[label=(\alph*), itemsep=6pt]
  \item $f(x) = -3(x-2)(x+5)^3$

        \begin{work}
          \text{degree of } f &= 1 + 3 \\
                              &= 4 \\
          \text{leading term} &= -3 \cdot x \cdot x^{3} \\
                              &= -3x^{4}
        \end{work}

        Degree: \ans{$4$} \quad
        $\lim_{x \to \infty} f(x) = $ \ans{$-\infty$} \quad
        $\lim_{x \to -\infty} f(x) = $ \ans{$-\infty$}

  \item $g(x) = 2x(x^2+7)(x-1)^2$

        \begin{work}
          \text{degree of } g &= 1 + 2 + 2 \\
                              &= 5 \\
          \text{leading term} &= 2 \cdot x \cdot x^{2} \cdot x^{2} \\
                              &= 2x^{5}
        \end{work}

        Degree: \ans{$5$} \quad
        $\lim_{x \to \infty} g(x) = $ \ans{$\infty$} \quad
        $\lim_{x \to -\infty} g(x) = $ \ans{$-\infty$}
\end{enumerate}

\item Dominance check for $p(x) = x^4 - 1000x^3$. The ratio of the rest to the
      leading term is $\dfrac{-1000x^3}{x^4} = \dfrac{-1000}{x}$.

\smallskip
\renewcommand{\arraystretch}{1.5}
\begin{center}
\begin{tabular}{|c|c|c|c|c|}
\hline
$x$ & $100$ & $1{,}000$ & $10{,}000$ & $100{,}000$ \\
\hline
rest $\div$ leading & \ans{$-10$} & \ans{$-1$} & \ans{$-0.1$} & \ans{$-0.01$} \\
\hline
\end{tabular}
\end{center}

\smallskip
So the $-1000x^3$ term is \ans{negligible} compared with $x^4$ once $x$ is large,
and therefore

$\lim_{x \to \infty} p(x) = $ \ans{$\infty$} \qquad
$\lim_{x \to -\infty} p(x) = $ \ans{$\infty$}

\boxguard[30]
\item Read the end behavior off each graph. You are not given either equation.

\begin{center}
\begin{tikzpicture}[x=0.65cm, y=0.26cm]
  \draw[linegray, very thin, xstep=1, ystep=2] (-2.4,-5) grid (2.4,5);
  \draw[->, thick] (-2.8,0) -- (2.9,0) node[below right] {\small $x$};
  \draw[->, thick] (0,-5.6) -- (0,5.6) node[above] {\small $y$};
  \draw[very thick, plum] plot [smooth] coordinates
    {(-2.2,4.07) (-2,0) (-1.7,-3.21) (-1.4,-4.0) (-1,-3) (-0.5,-0.94)
     (0,0) (0.5,-0.94) (1,-3) (1.4,-4.0) (1.7,-3.21) (2,0) (2.2,4.07)};
  \node[below] at (0,-6.2) {\small \textbf{Graph A}};
\end{tikzpicture}
\hspace{1.6cm}
\begin{tikzpicture}[x=0.65cm, y=0.26cm]
  \draw[linegray, very thin, xstep=1, ystep=2] (-2.4,-5) grid (2.4,5);
  \draw[->, thick] (-2.8,0) -- (2.9,0) node[below right] {\small $x$};
  \draw[->, thick] (0,-5.6) -- (0,5.6) node[above] {\small $y$};
  \draw[very thick, plum] plot [smooth] coordinates
    {(-2.2,-4.05) (-2,-2) (-1.8,-0.43) (-1.5,1.13) (-1,2) (-0.5,1.38)
     (0,0) (0.5,-1.38) (1,-2) (1.5,-1.13) (1.8,0.43) (2,2) (2.2,4.05)};
  \node[below] at (0,-6.2) {\small \textbf{Graph B}};
\end{tikzpicture}
\end{center}

\begin{enumerate}[label=(\alph*), itemsep=6pt]
  \item Graph A: \; $\lim_{x \to \infty} y = $ \ans{$\infty$}, \;
        $\lim_{x \to -\infty} y = $ \ans{$\infty$}, \;
        degree \ans{even}, \; $a_n$ \ans{positive}.
  \item Graph B: \; $\lim_{x \to \infty} y = $ \ans{$\infty$}, \;
        $\lim_{x \to -\infty} y = $ \ans{$-\infty$}, \;
        degree \ans{odd}, \; $a_n$ \ans{positive}.
\end{enumerate}

\item Which polynomial function satisfies both
      $\lim_{x \to -\infty} p(x) = \infty$ and
      $\lim_{x \to \infty} p(x) = -\infty$?

\begin{enumerate}[label=\Alph*., itemsep=3pt]
  \item $p(x) = 2x^4 - x$
  \item $p(x) = -2x^4 + x$
  \item $p(x) = 2x^3 - x$
  \item $p(x) = -2x^3 + x$ \quad \textcolor{keyred}{\textbf{$\leftarrow$ correct: odd degree, negative $a_n$}}
\end{enumerate}

\end{enumerate}

\vspace{0.1in}

\boxguard
\begin{extensionbox}
Factor the leading term out of a general polynomial:
\[
  p(x) = a_n x^n\left(1 + \frac{a_{n-1}}{a_n x} + \frac{a_{n-2}}{a_n x^2}
         + \dots + \frac{a_0}{a_n x^n}\right).
\]
Every fraction inside the parentheses has an $x$ in its denominator. Explain in
one or two sentences what the parentheses close in on as $x$ increases without
bound, and why that proves the dominance table from today's notes in one stroke.

\vspace{0.05in}
\ansline{Each fraction shrinks toward $0$, so the parentheses close in on $1$ and $p(x)$ behaves}
\ansline{like $a_n x^n$ --- the leading term --- no matter what the other coefficients are.}
\end{extensionbox}

\vspace{0.1in}

\boxguard
\begin{notesbox}{Preview of Lesson 2.6}
You now have both halves of a polynomial graph: Lesson 2.4 counted the zeros in
the middle, and today described the two far ends. Notice that each half came
from a different \textit{form} --- factored form handed over the zeros,
standard form handed over the end behavior. Lesson 2.6, \textit{Equivalent
Representations of Polynomial Expressions}, puts the two forms side by side and
shows how to move between them, so you can always reach for the one that
answers the question in front of you.
\end{notesbox}

\end{document}
